% =====================================================================
%  CARNET D'ENTRAÎNEMENT — FICHE 3 : Trigonométrie
%  Thème 3 — Relation fondamentale & formules d'addition/duplication
%  TUTORIEL
% =====================================================================
\documentclass[11pt,a4paper]{article}
\usepackage[utf8]{inputenc}
\usepackage[T1]{fontenc}
\usepackage{lmodern}
\usepackage[french]{babel}
\usepackage{amsmath,amssymb,mathtools}
\usepackage{icomma}
\usepackage{xcolor}
\usepackage{colortbl}
\usepackage{tikz}
\usepackage[most]{tcolorbox}
\usepackage{enumitem}
\usepackage{array}
\usepackage{booktabs}
\usepackage[a4paper,margin=2.2cm]{geometry}
\usetikzlibrary{arrows.meta,calc,angles,quotes,decorations.markings}

\definecolor{primary}{RGB}{214,51,108}
\definecolor{primarydark}{RGB}{140,20,64}
\definecolor{defcol}{RGB}{0,114,178}
\definecolor{excol}{RGB}{0,135,90}
\definecolor{methcol}{RGB}{90,90,100}
\definecolor{attncol}{RGB}{200,40,55}
\definecolor{sincol}{RGB}{205,40,55}
\definecolor{coscol}{RGB}{20,110,200}

\DeclareMathOperator{\tg}{tg}
\DeclareMathOperator{\cotg}{cotg}
\newcommand{\degr}{^{\circ}}
\newcommand{\Z}{\mathbb{Z}}

\tcbset{
  boxbase/.style={enhanced,breakable,boxrule=0.9pt,arc=2.5pt,
     left=3mm,right=3mm,top=2.3mm,bottom=2.3mm,
     fonttitle=\bfseries,coltitle=white,toptitle=0.6mm,bottomtitle=0.6mm},
}
\newtcolorbox{idee}[1][]{boxbase,colback=defcol!6,colframe=defcol,
  title={\textsc{L'essentiel}\ifx\relax#1\relax\relax\else\textmd{ : \itshape #1}\fi}}
\newtcolorbox{exemple}[1][]{boxbase,colback=excol!6,colframe=excol,
  title={\textsc{Exemple}\ifx\relax#1\relax\relax\else\textmd{ : \itshape #1}\fi}}
\newtcolorbox{methode}[1][]{boxbase,colback=methcol!8,colframe=methcol,sharp corners,
  title={\textsc{Méthode}\ifx\relax#1\relax\relax\else\textmd{ --- \itshape #1}\fi}}
\newtcolorbox{piege}[1][]{boxbase,colback=attncol!7,colframe=attncol,
  title={\textsc{Piège à éviter}\ifx\relax#1\relax\relax\else\textmd{ : \itshape #1}\fi}}

\newcommand{\heading}[1]{\medskip\noindent{\large\bfseries\color{primarydark}#1}\par\smallskip}

\pagestyle{empty}
\setlength{\parindent}{0pt}
\setlength{\parskip}{3pt}

\begin{document}

\begin{center}
{\LARGE\bfseries\color{primary}Thème 3 — Relation fondamentale \& formules d'addition}\\[2pt]
{\color{primarydark}\textsc{Carnet d'entraînement — Fiche 3 : Trigonométrie \quad•\quad Tutoriel}}
\end{center}
\hrule height 1pt
\medskip

Ce thème apprend à \textbf{calculer} : à partir d'\emph{un} nombre trigonométrique connu, retrouver les
autres (relation fondamentale) ; et à obtenir des valeurs exactes d'angles non remarquables en les
\textbf{décomposant} (formules d'addition et de duplication). Un rappel des valeurs remarquables :

\begin{center}
\renewcommand{\arraystretch}{1.5}
\begin{tabular}{|>{\columncolor{primary!12}}c|c|c|c|c|c|}
\hline
\rowcolor{primary!20}
$\alpha$ & $0$ & $\frac{\pi}{6}\,(30\degr)$ & $\frac{\pi}{4}\,(45\degr)$ & $\frac{\pi}{3}\,(60\degr)$ & $\frac{\pi}{2}\,(90\degr)$\\
\hline
$\sin\alpha$ & $0$ & $\frac{1}{2}$ & $\frac{\sqrt2}{2}$ & $\frac{\sqrt3}{2}$ & $1$\\ \hline
$\cos\alpha$ & $1$ & $\frac{\sqrt3}{2}$ & $\frac{\sqrt2}{2}$ & $\frac{1}{2}$ & $0$\\ \hline
$\tg\alpha$ & $0$ & $\frac{\sqrt3}{3}$ & $1$ & $\sqrt3$ & $\nexists$\\
\hline
\end{tabular}
\end{center}

\heading{A. La relation fondamentale}
Comme un point du cercle vérifie $x^2+y^2=1$ avec $x=\cos\alpha$, $y=\sin\alpha$ :
\[
\boxed{\;\cos^2\alpha+\sin^2\alpha=1\;}\qquad\Longrightarrow\qquad
\sin\alpha=\pm\sqrt{1-\cos^2\alpha},\qquad \cos\alpha=\pm\sqrt{1-\sin^2\alpha}.
\]
En divisant cette égalité par $\cos^2\alpha$, puis par $\sin^2\alpha$, on obtient deux relations pour la
tangente et la cotangente :
\[
1+\tg^2\alpha=\dfrac{1}{\cos^2\alpha},\qquad\qquad 1+\cotg^2\alpha=\dfrac{1}{\sin^2\alpha}.
\]

\begin{minipage}{0.60\textwidth}
\begin{idee}[le signe $\pm$ se règle avec le quadrant]
La racine carrée seule ne donne que la \emph{grandeur}. C'est le \textbf{quadrant} qui fixe le
\textbf{signe} : $\cos>0$ à droite (I, IV), $\sin>0$ en haut (I, II). Toujours l'écrire avant de
conclure --- c'est l'erreur la plus fréquente.
\end{idee}
\end{minipage}\hfill
\begin{minipage}{0.36\textwidth}
\centering
\begin{tikzpicture}[scale=1.0,line cap=round,>=Stealth]
  \coordinate (A) at (0,0);
  \coordinate (B) at (3.2,0);
  \coordinate (C) at (3.2,1.9);
  \draw[primary,line width=1pt] (A)--(B)--(C)--cycle;
  \draw[black!55] (2.9,0) |-(3.2,0.3);
  \node[below] at (1.6,-0.05) {\scriptsize adjacent};
  \node[right] at (3.2,0.95) {\scriptsize opposé};
  \node[above left] at (1.7,1.05) {\scriptsize hypoténuse};
  \pic[draw=primarydark,angle radius=8mm,"\scriptsize$\alpha$"]{angle=B--A--C};
\end{tikzpicture}\\[1mm]
{\scriptsize$\cos\alpha=\frac{\text{adj}}{\text{hyp}},\ \sin\alpha=\frac{\text{opp}}{\text{hyp}},\ \tg\alpha=\frac{\text{opp}}{\text{adj}}$}
\end{minipage}

\begin{methode}[un nombre connu $+$ le quadrant $\Rightarrow$ les autres]
\begin{enumerate}[leftmargin=1.6em,topsep=1pt,itemsep=1pt]
  \item Choisir la bonne relation : $\cos^2+\sin^2=1$ si on part de $\cos$ ou $\sin$ ;
        $1+\tg^2=\frac{1}{\cos^2}$ si on part de la tangente.
  \item Isoler le carré, prendre la racine : on obtient la \textbf{grandeur} ($\pm$).
  \item Fixer le \textbf{signe} d'après le quadrant donné.
  \item En déduire les autres par $\tg\alpha=\frac{\sin\alpha}{\cos\alpha}$.
\end{enumerate}
\end{methode}

\begin{exemple}[$\sin\alpha=-\tfrac15$, quadrant IV]
$\cos^2\alpha=1-\frac{1}{25}=\frac{24}{25}$. Au quadrant IV $\cos>0$ :
$\cos\alpha=+\frac{\sqrt{24}}{5}=\frac{2\sqrt6}{5}$. Puis
$\tg\alpha=\dfrac{-1/5}{2\sqrt6/5}=-\dfrac{1}{2\sqrt6}=-\dfrac{\sqrt6}{12}$.
\end{exemple}

\heading{B. Formules d'addition}
\[
\cos(a\pm b)=\cos a\cos b\mp\sin a\sin b,\qquad
\sin(a\pm b)=\sin a\cos b\pm\cos a\sin b,\qquad
\tg(a\pm b)=\dfrac{\tg a\pm\tg b}{1\mp\tg a\,\tg b}.
\]
Mémo des signes : pour le \textbf{cosinus} le signe est \emph{opposé} à celui de l'opération ; pour le
\textbf{sinus} il est \emph{identique}.

\begin{methode}[valeur exacte d'un angle non remarquable]
Écrire l'angle comme \textbf{somme ou différence} de deux angles remarquables
($15\degr=45\degr-30\degr$, $75\degr=45\degr+30\degr$, $105\degr=60\degr+45\degr$), puis appliquer la
formule et lire le tableau.
\end{methode}

\begin{exemple}[$\cos15\degr$]
$\cos15\degr=\cos(45\degr-30\degr)=\cos45\degr\cos30\degr+\sin45\degr\sin30\degr
=\frac{\sqrt2}{2}\cdot\frac{\sqrt3}{2}+\frac{\sqrt2}{2}\cdot\frac12=\dfrac{\sqrt6+\sqrt2}{4}.$
De même $\tg15\degr=\dfrac{\tg45\degr-\tg30\degr}{1+\tg45\degr\tg30\degr}=2-\sqrt3.$
\end{exemple}

\begin{piege}[le cosinus n'est pas linéaire]
$\cos(a+b)\neq\cos a+\cos b$ et $\sin(a+b)\neq\sin a+\sin b$. Ainsi
$\cos(60\degr+60\degr)=\cos120\degr=-\frac12$, alors que $\cos60\degr+\cos60\degr=1$ : totalement
différent. On \textbf{doit} passer par la formule.
\end{piege}

\heading{C. Formules de duplication (angle double)}
En posant $b=a$ dans les formules d'addition :
\[
\sin2a=2\sin a\cos a,\qquad
\cos2a=\cos^2a-\sin^2a=2\cos^2a-1=1-2\sin^2a,\qquad
\tg2a=\dfrac{2\tg a}{1-\tg^2a}.
\]
Lues à l'envers, les deux dernières donnent la \textbf{linéarisation} (baisser le carré) :
\[
\cos^2 x=\dfrac{1+\cos2x}{2},\qquad\qquad \sin^2 x=\dfrac{1-\cos2x}{2}.
\]

\begin{exemple}[chaîne complète : $\sin\alpha=\tfrac13$, $\alpha$ au quadrant I]
$\cos\alpha=+\sqrt{1-\frac19}=\frac{2\sqrt2}{3}$. Puis
$\sin2\alpha=2\cdot\frac13\cdot\frac{2\sqrt2}{3}=\frac{4\sqrt2}{9}$ et
$\cos2\alpha=1-2\cdot\frac19=\frac79$. Enfin, par addition,
\[
\sin\!\Big(2\alpha-\tfrac{\pi}{3}\Big)=\sin2\alpha\cos\tfrac{\pi}{3}-\cos2\alpha\sin\tfrac{\pi}{3}
=\frac{4\sqrt2}{9}\cdot\frac12-\frac79\cdot\frac{\sqrt3}{2}=\dfrac{4\sqrt2-7\sqrt3}{18}.
\]
\end{exemple}

\end{document}
